\textbf{Время:} 01:45

\textbf{Тезис:} корреляционная функция получается усреднением интерференционного веса по точкам испускания.

\textbf{На экране:} источник и две выделенные точки слева; первая формула — среднее веса; вторая — запись через распределение ρ; ниже — предположения.

\textbf{Текст:} На предыдущем слайде были две фиксированные точки. Слева теперь показана вся область испускания: выбираем в ней разные A и B. Скобки усреднения в первой формуле означают, что складываем вклады всех таких пар с их вероятностями. Единица остаётся единицей, а косинус заменяется средним косинусом. Так появляется корреляционная функция C.

Вероятность выбрать точку задаёт распределение rho. Мы рассматриваем пары с определённым средним импульсом K, поэтому у rho есть этот аргумент. Распределение нормировано на единицу.

Вторая формула записывает то же среднее через интеграл по источнику. Здесь уже узнаваемая связь с волнами: вклад каждой точки имеет свою фазу. При малой разности импульсов фазы близки; при большей начинают взаимно компенсироваться. Скорость этого изменения зависит от пространственных и временных разделений точек.

Эта связь получена при независимом испускании и без взаимодействия пионов после испускания — предположения подписаны под формулой. На следующем примере посмотрим, как размер меняет ширину пика.

\textbf{Переход:} Выберем простой источник, для которого размер заранее известен, и сравним две корреляционные кривые.

\textbf{Справка:} хаотическое независимое испускание, приближение гладкости по импульсу, плоские волны, отсутствие многобозонных эффектов. Промежуточная запись для вопросов:

\[\int\rho(x;K)\,\dd^4x=1,\qquad\rho_A=\rho(x_A;K),\quad\rho_B=\rho(x_B;K).\]

\[C-1\simeq\iint\dd^4x_A\,\dd^4x_B\,\rho_A\rho_B\cos\!\frac{q\cdot(x_A-x_B)}{\hc}.\]

Экспонента разности фаз раскладывается на множитель для A и сопряжённый множитель для B. Независимые интегралы дают квадрат модуля одного интеграла, показанный на экране. Потеря фазы и ограничение массовой оболочки не позволяют однозначно восстановить произвольный четырёхмерный источник.

\textbf{Источники:} Lisa et al. (2005), §2.1–2.2.
